% Einheit 4 von 4 – Sinussatz (bank/trigonometrie/e4.jsonl, zusammenbau v0.1)
%% TODO Zweigzeile Teil 1: was hier gelernt wird (Satz in Schülersprache aus dem Einheitstitel) – Einheit 4
\einheitenkopf[e4]{Einheit 4 von 4 · Sinussatz}
\zweigzeile{ab Kl. 10 (am Gymnasium ab Kl. 9) · P10 oft}
\setcounter{aufgabe}{25}

%% TODO Titel: Ich-kann-Satz fehlt in der Bank (Platzhalter: Kettenname) – Nr. 26
%% TODO Anweisung: ein Satz über der ersten Teilaufgabe fehlt in der Bank – Nr. 26
\begin{aufgabe}{Sinussatz}
%% TODO \rechenplatz{n}: Schrittzahl des Grundfalls fehlt in der Bank (nur bei fester Schreibform, 2.3 b)
\begin{teile}
\teil Im Dreieck $ABC$ ist bei $B$ ein rechter Winkel markiert, $a$ und $\alpha$ sind gegeben. Womit rechnest du $b$? \\ \kreuz{sin, cos, tan} \kreuz{Sinussatz}
\teil Dreieck $ABC$: $a = 7$ cm, $\alpha = 68^\circ$, $\beta = 51^\circ$. Wie lang ist $b$? b ≈ \leerfeld[cm]
\teil Dreieck $PQR$: $PQ = 8$ cm, Winkel bei $R$: $69^\circ$, Winkel bei $Q$: $53^\circ$. Wie lang ist $PR$? PR ≈ \leerfeld[cm]
\teil Dreieck $ABC$: $a = 8$ cm, $\alpha = 63^\circ$, $\beta = 48^\circ$. Wie lang ist $c$? c ≈ \leerfeld[cm]
\teil Im Drachen $ABCD$ gilt im Dreieck $ABD$: $AD = 4{,}6$ cm, der Winkel bei $A$ ist $118^\circ$, der Winkel bei $B$ ist $33^\circ$. Berechne die Diagonale $BD$. \hfill (P10 2015 OS) \\ BD ≈ \leerfeld[cm]
\teil Dreieck $ABC$: $a = 7$ cm, $\alpha = 27^\circ$, $\beta = 118^\circ$. Wie lang ist $c$? c ≈ \leerfeld[cm]
\teil Dreieck $ABC$: $a = 4{,}35$ m, $\alpha = 58{,}2^\circ$, $\beta = 47{,}6^\circ$. Wie lang ist $b$, auf eine Dezimale? b ≈ \leerfeld[m]
\teil Dreieck $ABC$: $a = 9$ cm, $\alpha = 77^\circ$, $\beta = 44^\circ$. Zeige rechnerisch, dass $b \approx 6{,}4$ cm ist.
\teil Dreieck $ABC$: $c = 10$ cm, $\alpha = 53^\circ$, $\beta = 56^\circ$. Prüfe die Aussage „$a$ ist länger als $b$.“
\teil Viereck $ABCD$ mit rechtem Winkel bei $B$ und der Diagonale $AC = 8$ m. $\angle BAC = 57^\circ$, $\angle BCD = 91^\circ$, $\angle ADC = 76^\circ$. Wie lang ist $AD$? AD ≈ \leerfeld[m]
\teil Wanderung: Hütte $H$, Gipfel $G$, See $S$. $HG = 3\,200$ m, der Winkel bei $H$ ist $57^\circ$, der Winkel bei $S$ ist $72^\circ$. Wie lang ist der Weg vom Gipfel zum See und weiter zum Parkplatz, der $1\,800$ m hinter dem See liegt? \hfill (P10 2014 OS) \\ \leerfeld[km]
\teil Dreieck $ABC$: $a = 8$ cm, $\alpha = 67^\circ$, $\beta = 53^\circ$. Berechne $b$ mit dem Sinussatz und noch einmal über die Höhe $h$ von $C$ auf $AB$. b ≈ \leerfeld[cm]
\teil Dreieck $ABC$: $a = 9$ cm, $\alpha = 68^\circ$, $b = 7$ cm. Wie groß ist $\beta$? $\beta \approx$ \leerfeld $^\circ$
\teil Dreieck $ABC$: $a = 5$ cm, $b = 8$ cm, $\gamma = 47^\circ$. Wie lang ist $c$? \leerfeld[cm]
\teil Dreieck $ABC$: $a = 5$ cm, $b = 7$ cm, $c = 6$ cm. Wie groß ist $\gamma$? $\gamma \approx$ \leerfeld $^\circ$
\teil Die Höhe $h_c$ zerlegt das Dreieck $ABC$ in zwei rechtwinklige Teildreiecke. Drücke $h_c$ in beiden Teildreiecken aus und leite daraus $\frac{a}{\mathrm{sin}\,\alpha} = \frac{b}{\mathrm{sin}\,\beta}$ her.
\teil Eine Seilbahn führt von $A$ über die Stütze $B$ nach $C$. $AB = 520$ m, der Winkel bei $A$ ist $31^\circ$, der Winkel $ABC$ ist $113^\circ$. Berechne die Länge der Strecke $BC$. \hfill (P10 2024 OS) \\ BC ≈ \leerfeld[m]
\end{teile}
\end{aufgabe}

%% TODO Titel: Ich-kann-Satz fehlt in der Bank (Platzhalter: Kettenname) – Nr. 27
%% TODO Anweisung: ein Satz über der ersten Teilaufgabe fehlt in der Bank – Nr. 27
\begin{aufgabe}{Sinussatz – Fehler finden, Begründen, Anwendung}
\begin{teile}
\teil Dreieck $ABC$: $a = 8$ cm, $\alpha = 71^\circ$, $\beta = 43^\circ$, gesucht $b$. Wo ist der Fehler? \rechnung{b &= 8 \cdot \mathrm{sin}\,71^\circ : \mathrm{sin}\,43^\circ \\ b &\approx 11{,}1 \text{ cm}}
\teil Die Höhe $h$ von $C$ auf $AB$ ist in beiden Teildreiecken dieselbe Kathete. Warum folgt daraus $b \cdot \mathrm{sin}\,\alpha = a \cdot \mathrm{sin}\,\beta$?
\teil Über einen Fluss soll ein Seil von $A$ zu einem Baum $C$ am anderen Ufer gespannt werden. Am eigenen Ufer ist $AB = 80$ m abgesteckt, der Winkel bei $A$ ist $72^\circ$, der bei $B$ $63^\circ$. Reicht eine $100$ m lange Seilrolle?
\end{teile}
\end{aufgabe}
